\documentclass[../Odyssey_Project.tex]{subfiles}

\begin{document}
\setcounter{section}{0}
\section{Review}
\subsubsection*{Group Axioms}
A group consists of a non-empty set G and a binary operation on G such that the following axioms hold:
\begin{itemize}
    \item \(\forall g, h, k \in G, (gh)k = g(hk)\) (associativity)
    \item \(\exists e \in G\) s.t. \(eg = ge = g \ \forall g \in G\) (identity)
    \item \(\forall g \in G, \exists g^{-1} \in G\) s.t. \(gg^{-1} = g^{-1}g = e\) (inverses)
\end{itemize}
\subsubsection*{Properties of Groups}
\begin{itemize}
    \item If \(xy = yx\), the commutator \([x,y] = xyx^{-1}y^{-1} = e\)
    \item If \(xy = yx \ \forall x,y \in G \), then G is an \textbf{abelian} group. 
    \item A \(x \in G\) is of finite order if \(\exists n \in \mathbb{N}\) s.t. \(x^n = e\)
    \item \(|G|\), the number of elements in G is defined to be the order of a finite group G. Every element of a finite group is of finite order
    \item There are infinite groups that has every element having a finite order, but there are also infinite groups in which the identity is the only emenet of finite order (torsion-free groups such as \(\mathbb{Z}\))
\end{itemize}
\subsubsection*{Properties of Subgroups}
\begin{itemize}
    \item H is a \textbf{subgroup} of G, denoted \(H \leq G\), if \(H\subseteq G\) and H is a group under the same operation as G. In particular, \(e \in H\) and \(h^{-1} \in H \ \forall h \in H\)
    \item Every subgroup of a finite group is finite, but an infinite group always has both finite and infinite subgroups
    \item Every subgroup of an abelian group is abelian, but a non-abelian group can have both abelian and non-abelian subgroups
\end{itemize}
\begin{proposition}
\label{prop:1.1}
If \(H,K \leq G\), then $H \cap K \leq G$. More generally, \(H_i\leq G \Rightarrow \bigcap H_i \leq G\)
\end{proposition}
\begin{proof}
    Suppose \(H_i \leq G \ \forall i\) \\
    Since each \(H_i \leq G\), we have \(e \in H_i \ \forall i\), so \(e \in \bigcap H_i\) and \(\bigcap H_i \neq \emptyset\). For \(x, y \in \bigcap H_i\), we have \(x, y \in H_i \ \forall i\), so by closure of each \(H_i\), \(xy \in H_i \ \forall i\), giving \(xy \in \bigcap H_i\). Similarly, \(x^{-1} \in H_i \ \forall i\), so \(x^{-1} \in \bigcap H_i\). \\
    \(\therefore \bigcap H_i \leq G\)
\end{proof}
\begin{namedtheorem}[Lagrange's Theorem]
\label{thm:lagrange}
Let G be a \textbf{finite} group, and let $H \leq G$. Then $|H| \mid |G|$.
\end{namedtheorem}
\begin{proof}
    Let \(x \in G\). Define \(f_x \colon H \to xH\) by \(f_x(h) = xh\). If \(f_x(h_1) = f_x(h_2)\), then \(xh_1 = xh_2\), and multiplying on the left by \(x^{-1}\) gives \(h_1 = h_2\). Also, every element of \(xH\) has the form \(xh\) for some \(h \in H\), so \(f_x\) is surjective. Thus \(f_x\) is a bijection, and \(|xH| = |H|\). \\
    We claim that any two sets of the form \(xH\) are either equal or disjoint. Suppose that \(xH \cap yH \neq \emptyset\). Then \(xh_1 = yh_2\) for some \(h_1, h_2 \in H\), so \(x = y(h_2h_1^{-1})\). Since \(h_2h_1^{-1} \in H\), we have \(x = yh\) for some \(h \in H\), and so \(xH = yhH = yH\). Thus the distinct sets of the form \(xH\) partition \(G\). \\
    Since \(G\) is finite, there are finitely many such sets, say \(r\). Each has cardinality \(|H|\), so \(|G| = r|H|\).  \\
    \(\therefore |H| \mid |G|\).
\end{proof}
\smallskip
\noindent If \(X \subseteq G, \langle X \rangle \coloneqq \bigcap \{ H \leq G \mid X \subseteq H \}\). By Proposition~\ref{prop:1.1}, \(\langle X \rangle \leq G\), and it is the smallest subgroup of G containing X. We call \(\langle X \rangle\) the \textbf{subgroup generated by X}, and we say that X \textbf{generates} \(\langle X \rangle\).
\smallskip
\begin{proposition}
\label{prop:1.2}
Let \(X \subseteq G\). Then \(\langle X \rangle = \{e\} \cup \{x_1^{\epsilon_1} \cdots x_r^{\epsilon_r} \mid r \geq 1,\ x_i \in X,\ \epsilon_i = \pm 1,\ 1 \leq i \leq r\}\).
\end{proposition}
\begin{proof}
    Let \(S = \{e\} \cup \{x_1^{\epsilon_1} \cdots x_r^{\epsilon_r} \mid r \geq 1,\ x_i \in X,\ \epsilon_i = \pm 1,\ 1 \leq i \leq r\}\). \\
    We first show that \(S \leq G\). Clearly \(e \in S\).
    If \(a = e\) or \(b = e\), then \(ab \in S\).
    Otherwise, if \(a = x_1^{\epsilon_1} \cdots x_r^{\epsilon_r} \in S\) and \(b = y_1^{\delta_1} \cdots y_s^{\delta_s} \in S\), then \(ab = x_1^{\epsilon_1} \cdots x_r^{\epsilon_r} y_1^{\delta_1} \cdots y_s^{\delta_s} \in S\).
    Also \(e^{-1} = e\), and if \(a = x_1^{\epsilon_1} \cdots x_r^{\epsilon_r}\), then \(a^{-1} = x_r^{-\epsilon_r} \cdots x_1^{-\epsilon_1} \in S\).
    Hence \(S \leq G\). \\
    Let \(\mathcal{H}_X = \{H \leq G \mid X \subseteq H\}\).
    By Proposition~\ref{prop:1.1}, \(\langle X \rangle = \bigcap_{H \in \mathcal{H}_X} H \leq G\). 
    Since \(X \subseteq S\), we have \(S \in \mathcal{H}_X\), and so \(\langle X \rangle = \bigcap_{H \in \mathcal{H}_X} H \subseteq S\). \\
    Conversely, let \(H \in \mathcal{H}_X\).
    Since \(x_i \in X \subseteq H\), we have \(x_i^{\epsilon_i} \in H\), so \(x_1^{\epsilon_1} \cdots x_r^{\epsilon_r} \in H\).
    Also \(e \in H\), hence \(S \subseteq H\).
    Thus \(S \subseteq H\) for all \(H \in \mathcal{H}_X\), so \(S \subseteq \bigcap_{H \in \mathcal{H}_X} H = \langle X \rangle\). \\
    \(\therefore S = \langle X \rangle\).
\end{proof}
\subsubsection*{Cyclic Groups}
\begin{itemize}
    \item A group G is \textbf{cyclic} if \(G = \langle g \rangle\) for some \(g \in G\). By Proposition~\ref{prop:1.2}, \(\langle g \rangle = \{g^n \mid n \in \mathbb{Z}\}\) s.t. \(|\langle g \rangle | = n\) and \(o(g) = n\)
    \item All cyclic groups are abelian
    \item If g is not of finite order, then \(\langle g \rangle \) is a torsion-free infinite abelian group, and \(\langle g \rangle \cong \Z\)
    \item If g is of finite order, then \(\langle g \rangle \cong \Z / n \Z\)
    \item If \(g \in G\) and \(o(g) = n \in \Z\), \(\langle g \rangle \leq G\), so by \nameref{thm:lagrange}, \(n \mid |G|\). Thus, the order of an element of a finite group must divide the order of that group
\end{itemize}
If \(X,Y \subseteq G, XY \coloneqq \{xy \mid x \in X, y \in Y\} \subseteq G \). Also, \(X^{-1} \coloneqq \{x^{-1} \mid x \in X\} \subseteq G\). If \(H \ne \emptyset\), then \(H \leq G \iff HH = H\) and \(H^{-1} = H\)
\begin{proposition}
\label{prop:1.3}
Let \(H,K \leq G\). Then \(HK \leq G \iff HK = KH\)
\end{proposition}
\begin{proof}
    \((\Rightarrow)\) Suppose \(HK \leq G\). Since \(HK\) is a subgroup, \((HK)^{-1} = HK\). Also \(H^{-1} = H\) and \(K^{-1} = K\), since \(H,K \leq G\). Hence \((HK)^{-1} = K^{-1}H^{-1} = KH\). Therefore \(HK = (HK)^{-1} = KH\). \\
    \((\Leftarrow)\) Suppose \(HK = KH\). Since \(e \in H \cap K\), we have \(e = ee \in HK\). Let \(h_1k_1, h_2k_2 \in HK\). Since \(k_1h_2 \in KH = HK\), write \(k_1h_2 = h_3k_3\) for some \(h_3 \in H,\ k_3 \in K\). Then \((h_1k_1)(h_2k_2) = h_1(k_1h_2)k_2 = h_1h_3k_3k_2 \in HK\). Also \((hk)^{-1} = k^{-1}h^{-1} \in KH = HK\) for all \(h \in H,\ k \in K\). \\
    \(\therefore HK \leq G\).
\end{proof}

\begin{theorem}
\label{thm:1.4}
Let G = \(\langle g \rangle \) be a cyclic group of order n
\begin{itemize}
    \item For each \(d \mid n\), there is exactly one subgroup of order d, \(\langle g^{\frac{n}{d}}\rangle\)
    \item If \(d,e \mid n\), then \(\langle g^\frac{n}{d} \rangle \cap \langle g^\frac{n}{e}\rangle = \langle g^\frac{n}{\gcd(d,e)}\rangle\)
    \item If \(d,e \mid n\), then \(\langle g^\frac{n}{d} \rangle \langle g^\frac{n}{e}\rangle = \langle g^\frac{n}{\operatorname{lcm}(d,e)}\rangle\)
\end{itemize}
\end{theorem}
\begin{proof}
    First note that if \(m \mid n\), then \(o(g^m) = \frac{n}{m}\). Indeed, \((g^m)^{\frac{n}{m}} = g^n = e\), and if \((g^m)^t = e\), then \(n \mid mt\), so \(\frac{n}{m} \mid t\). \\
    Let \(d \mid n\). Taking \(m = \frac{n}{d}\), we get \(o(g^{\frac{n}{d}}) = d\), so \(\langle g^{\frac{n}{d}}\rangle\) is a subgroup of order \(d\). \\
    For uniqueness, let \(H \leq G\) with \(|H| = d\), and let \(a = \min\{m \in \N \mid g^m \in H\}\). This set is non-empty since \(g^n = e \in H\). If \(g^b \in H\), write \(b = qa + r\), \(0 \leq r < a\). Then \(g^r = g^b(g^a)^{-q} \in H\), so by minimality of \(a\), \(r = 0\). Hence \(a \mid b\), so \(H \subseteq \langle g^a\rangle\). Since \(g^a \in H\), \(\langle g^a\rangle \subseteq H\), and hence \(H = \langle g^a\rangle\). Applying the same division argument to \(n = qa+r\) gives \(a \mid n\), so \(|H| = |\langle g^a\rangle| = \frac{n}{a} = d\). Thus \(a = \frac{n}{d}\), and \(H = \langle g^{\frac{n}{d}}\rangle\). \\
    Now let \(d,e \mid n\). An element \(g^b \in \langle g^{\frac{n}{d}}\rangle \cap \langle g^{\frac{n}{e}}\rangle\) iff \(\frac{n}{d} \mid b\) and \(\frac{n}{e} \mid b\), iff \(\operatorname{lcm}(\frac{n}{d},\frac{n}{e}) = \frac{n}{\gcd(d,e)}\) divides \(b\). Hence \(\langle g^{\frac{n}{d}}\rangle \cap \langle g^{\frac{n}{e}}\rangle = \langle g^{\frac{n}{\gcd(d,e)}}\rangle\). \\
    Finally, \(\langle g^{\frac{n}{d}}\rangle\langle g^{\frac{n}{e}}\rangle = \{g^{\frac{n}{d}r+\frac{n}{e}s} \mid r,s \in \Z\}\). Since \(\frac{n}{d}\Z + \frac{n}{e}\Z = \gcd(\frac{n}{d},\frac{n}{e})\Z = \frac{n}{\operatorname{lcm}(d,e)}\Z\), we get \(\langle g^{\frac{n}{d}}\rangle\langle g^{\frac{n}{e}}\rangle = \langle g^{\frac{n}{\operatorname{lcm}(d,e)}}\rangle\). \\
    \(\therefore\) all three assertions hold.
\end{proof}
\subsubsection*{Cosets}
\begin{itemize}
    \item If \(H \leq G\) and \(x \in G, xH\) is called a left coset of H in G and \(Hx\) is called a right coset of H in g
    \item There is a bijective correspondence between left and right cosets of H in G, sending \(xH\) to its inverse \((xH)^{-1} = Hx^{-1}\)
    \item Any two cosets of are equal or disjoint, with \(xH = yH \iff y^{-1}x \in H\)
    \item A \(x \in G\) lies in exacty one coset of H, \(xH\)
    \item \(\forall x \in G, \exists \) a bijective correspondence between H and xH
    \item We define the index of H in G, \(|G:H|\) to be the number of cosets of H in G
    \item The cosets of H in G partition G into \(|G:H|\) disjoint sets of cardinality \(|H|\), and we have \(|G| = |G:H||H|\) (\nameref{thm:lagrange})
    \item All subgroups of a finite group are of finite index, while subgroups of an infinite group may be of finite or infinite index
\end{itemize}

\begin{theorem}
\label{thm:1.5}
Let \(G = \langle g \rangle \) be an infinite cyclic group
\begin{itemize}
    \item For each \(d \in \N, \exists \) exactly one subgroup of G of index d, \(\langle g^d \rangle \). Furthermore, every non-trivial subgroup of G is of finite index
    \item The intersection of the subgroups of indices d and e is the subgroup of index \(\operatorname{lcm}(d,e)\)
    \item The product of subgroups of indices d and e is the subgroup of index \(\gcd(d,e)\)
\end{itemize}
\end{theorem}
\begin{proof}
    As in the proof of Theorem~\ref{thm:1.4}, let \(1 \ne H \leq G\), and choose \(a = \min\{m \in \N \mid g^m \in H\}\). This set is non-empty since \(H\) contains some \(g^b \ne e\), and then either \(b > 0\) or \(g^{-b} \in H\) with \(-b > 0\). If \(g^c \in H\), write \(c = qa+r\), \(0 \leq r < a\). Then \(g^r = g^c(g^a)^{-q} \in H\), so minimality of \(a\) gives \(r = 0\). Hence \(a \mid c\), so \(H = \langle g^a\rangle\). \\
    For \(d \in \N\), the cosets of \(\langle g^d\rangle\) are \(g^i\langle g^d\rangle\), \(0 \leq i < d\). Hence \(|G:\langle g^d\rangle| = d\). By the previous paragraph, every non-trivial subgroup is \(\langle g^a\rangle\) for some \(a \in \N\), so it has finite index \(a\), and \(\langle g^d\rangle\) is the unique subgroup of index \(d\). \\
    If \(d,e \in \N\), then \(g^b \in \langle g^d\rangle \cap \langle g^e\rangle \iff d \mid b\) and \(e \mid b \iff \operatorname{lcm}(d,e) \mid b\). Thus \(\langle g^d\rangle \cap \langle g^e\rangle = \langle g^{\operatorname{lcm}(d,e)}\rangle\), which has index \(\operatorname{lcm}(d,e)\). \\
    Since \(G\) is abelian, \(\langle g^d\rangle\langle g^e\rangle \leq G\) by Proposition~\ref{prop:1.3}. Moreover \(\langle g^d\rangle\langle g^e\rangle = \{g^{dr+es} \mid r,s \in \Z\} = \langle g^{\gcd(d,e)}\rangle\), since \(d\Z + e\Z = \gcd(d,e)\Z\). Hence the product has index \(\gcd(d,e)\). \\
    \(\therefore\) all three assertions hold.
\end{proof}

\begin{theorem}
\label{thm:1.6}
If \(K \leq H \leq G\), then \(|G:K| = |G:H| |H:K|\)
\end{theorem}
\begin{proof}
    By the coset-counting identity in \nameref{thm:lagrange}, the coefficient of \(|H|\) in \(|G|\) is \(|G:H|\), and the coefficient of \(|K|\) in \(|H|\) is \(|H:K|\). Hence \(|G| = |G:H||H|\) and \(|H| = |H:K||K|\). \\
    Substituting, \(|G| = |G:H||H:K||K|\). But applying the same identity directly to \(K \leq G\) gives \(|G| = |G:K||K|\). \\
    \(\therefore |G:K| = |G:H||H:K|\).
\end{proof}

\smallskip
\noindent Let $H \leq G$, and let $\mathcal{I}$ be an indexing set that is in bijective correspondence with $G/H$. $T = \{t_i \mid i \in \mathcal{I}\} \subseteq G$ is said to be a (\emph{left}) \emph{transversal} for $H$ (or a set of (\emph{left}) \emph{coset representatives} of $H$ in $G$) if the sets $t_i H$ are precisely the distinct cosets of $H$ in $G$
\subsubsection*{Normal Subgroups}
\begin{itemize}
    \item $N$ is a \emph{normal subgroup} of $G$ if $xN = N x \ \forall x \in G$, or equivalently if $xNx^{-1} \subseteq N$ for all $x \in G$
    \item If $G$ is abelian, then every subgroup of $G$ is normal. 
    \item If \{e\} and G are the only normal subgroups of $G$, then we say that $G$ is \emph{simple}. 
    \item If $H \trianglelefteq G$ and $K \trianglelefteq H$, then it is not necessarily true that $K \trianglelefteq G$. However, it is clearly true that if $K \trianglelefteq G$ and $K \leq H \leq G$, then $K \trianglelefteq H$.
\end{itemize}

\begin{proposition}
\label{prop:1.7}
Let $H, K \leq G$. If $K \trianglelefteq G$, then $HK \leq G$ and $H \cap K \trianglelefteq H$; if also $H \trianglelefteq G$, then $HK \trianglelefteq G$ and $H \cap K \trianglelefteq G$.
\end{proposition}
\begin{proof}
    Let \(H, K \leq G\). Suppose that \(K \trianglelefteq G\). \\
    Then we have \(gK = Kg \ \forall g \in G\). Since \(H \leq G\), we have \(hK = Kh \ \forall h \in H\). \\
    Take \(x \in HK,\) so \(x = hk\) for some \(h \in H, k \in K\). Since \(hK = Kh, \exists\) some \(h' \in H\) such that \(x = kh'\), so we have \(HK \subseteq KH\). By a similar argument, we have \(KH \subseteq HK\), so we conclude \(HK = KH\). By Proposition~\ref{prop:1.3}, we have \(HK \leq G\). \\
    Take \(y \in H \cap K\) and \(h \in H\). Since \(y \in H\), we have \(hyh^{-1} \in H\). Since \(y \in K\) and \(K \trianglelefteq G\), we have \(hKh^{-1} \subseteq K \Rightarrow hyh^{-1} \in K\). It follows that \(hyh^{-1} \in H \cap K\), so \(H \cap K \trianglelefteq H\). \\
    Suppose also that \(H \trianglelefteq G\). Take \(hk \in HK\) and \(g \in G\). Then \(g(hk)g^{-1} = (ghg^{-1})(gkg^{-1})\), where \(ghg^{-1} \in H\) since \(H \trianglelefteq G\) and \(gkg^{-1} \in K\) since \(K \trianglelefteq G\). Hence \(g(hk)g^{-1} \in HK\), so \(gHKg^{-1} \subseteq HK \Rightarrow HK \trianglelefteq G\). \\
    Take \(y \in H \cap K\) and \(g \in G\). Since \(H \trianglelefteq G\), \(gyg^{-1} \in H\), and since \(K \trianglelefteq G\), \(gyg^{-1} \in K\). So \(gyg^{-1} \in H \cap K \Rightarrow H \cap K \trianglelefteq G\).
\end{proof}
\begin{proposition}
\label{prop:1.8}
Any subgroup of index 2 is normal
\end{proposition}
\begin{proof}
    Let \(H \leq G\) and suppose that \(|G:H| = 2\) \\
    Then there are 2 left cosets, H and G - H \\
    It follows that \(x \in H \iff xH = H = Hx, \) and \(x \not\in H \iff xH = G - H = Hx\) \\
    \(\therefore H \trianglelefteq G\)
\end{proof}
\begin{theorem}
\label{thm:1.9}
If \(N \trianglelefteq G, \) then \(G / N\) forms the quotient group under the operation \((xN)(yN) = (xy)N\)
\end{theorem}
\begin{proof}
    First, we show that the operation is well-defined. Suppose \(xN = x'N\) and \(yN = y'N\). Then \(x' = xn_1\) and \(y' = yn_2\) for some \(n_1,n_2 \in N\). Since \(N \trianglelefteq G\), \(y^{-1}n_1y \in N\), so \(n_1y = y(y^{-1}n_1y)\). Hence \(x'y' = xn_1yn_2 = xy(y^{-1}n_1y)n_2 \in xyN\), and so \(x'y'N = xyN\). Thus \((xN)(yN) = (xy)N\) is well-defined. \\
    Associativity follows from associativity in \(G\), since \(((xN)(yN))(zN) = ((xy)z)N\) and \((xN)((yN)(zN)) = (x(yz))N\). \\
    The identity is \(N = eN\), since \(N(xN) = xN = (xN)N\). Also, \((xN)(x^{-1}N) = N = (x^{-1}N)(xN)\), so \(x^{-1}N\) is the inverse of \(xN\). \\
    \(\therefore G/N\) is a group under the operation \((xN)(yN) = (xy)N\).
\end{proof}
\smallskip
\noindent The identity element of \(G/N\) is \(N\), and the inverse is \(x^{-1}N\). If G is abelian, then \(G/N\) is also abelian
\subsubsection*{Conjugation}
\begin{itemize}
    \item The conjugate of an element \(x\) by \(g\) is defined as \(gxg^{-1}\)
    \item \(x\) and \(y\) are conjugate if \(\exists g \in G\) s.t. \(y = gxg^{-1}\)
    \item No 2 distinct elements of an abelian group can be conjugate
    \item \(N \trianglelefteq G \iff \ \forall n \in N, gng^{-1} \in N\)
\end{itemize}
\subsubsection*{Permutations}
\begin{itemize}
    \item A permutation of a set X is a bijective set map from X to X. The set of permutations of X forms a group under composition
    \item If \(X = \{1,\dots , n\}\), then this groups is called the symmetric group, \(S_n\)
    \item \(S_n\) is a finite group of order \(n!\)
    \item \(\rho \in S_n\) is called a r-cycle if \(\exists\) distinct integers \(1 \leq a_1,\dots, a_r \leq n\) s.t. \(\rho(a_i) = a_{i+1} \forall \ 1 \leq i \le r, \rho(a_r) = a_1, \rho(b) = b \ \forall 1 \leq b \leq n\) which is not equal to some \(a_i\)
    \item Two cycles are disjoint if there is no number that is moved by both cycles
    \item Every element of \(S_n\) can be written as a product of disjoint cycles
    \item The cycle structure refers to the partion of \(n\) consisting of the lengths of the cycles that appear in a disjoint cycle decomposition of \(\rho\)
\end{itemize}
\begin{proposition}
\label{prop:1.10}
Let \(n \in \N\). Then two elements of \(S_n\) are conjugate iff they have the same cycle structure
\end{proposition}
\begin{proof}
    Let \(\alpha, \beta \in S_n\). \\
    (\(\Rightarrow\)) Suppose that \(\alpha\) and \(\beta\) are conjugate, so \(\beta = \sigma \alpha \sigma^{-1}\) for some \(\sigma \in S_n\). We show that \(\sigma \alpha \sigma^{-1}\) is the permutation with the same cycle structure as \(\alpha\) obtained by applying \(\sigma\) to the symbols in \(\alpha\). Let \(\pi\) be this permutation. \\
    Suppose that \(\alpha\) fixes some \(i\), so \(\sigma(i)\) resides in a 1-cycle of \(\pi\), and \(\pi\) fixes \(\sigma(i)\). Then \(\sigma \alpha \sigma^{-1}(\sigma(i)) = \sigma\alpha(i) = \sigma(i)\), so \(\sigma \alpha \sigma^{-1}\) fixes \(\sigma(i)\) as well. \\
    Suppose that \(\alpha\) moves \(i\), and let \(\alpha(i) = j\). Take the complete factorization of \(\alpha\) to be
    \[\alpha = \gamma_1 \gamma_2 \cdots (\cdots i \ j \cdots) \cdots \gamma_t\]
    Take \(\sigma(i) = k\) and \(\sigma(j) = l\), so \(\pi: k \mapsto l\). Since \(\sigma\alpha\sigma^{-1}: k \mapsto i \mapsto j \mapsto l\), we have \(\sigma\alpha\sigma^{-1}(k) = \pi(k)\). So \(\pi\) and \(\sigma\alpha\sigma^{-1}\) agree on all symbols of the form \(k = \sigma(i)\). Since \(\sigma\) is a surjection, it follows that \(\pi = \sigma\alpha\sigma^{-1} = \beta\), so \(\alpha\) and \(\beta\) have the same cycle structure. \\
    (\(\Leftarrow\)) Suppose that \(\alpha\) and \(\beta\) have the same cycle structure. Place the complete factorization of \(\alpha\) over that of \(\beta\) so that cycles of the same length correspond, and let \(\gamma \in S_n\) be the function sending the top to the bottom. For example, if
    \[\alpha = \gamma_1 \gamma_2 \cdots (\cdots i \ j \cdots) \cdots \gamma_t\]
    \[\beta = \delta_1 \delta_2 \cdots (\cdots k \ l \cdots) \cdots \delta_t\]
	    then \(\gamma(i) = k, \gamma(j) = l\), etc. \(\gamma\) is a permutation, since every \(i\) between 1 and \(n\) occurs exactly once in a complete factorization. By the forward direction, \(\gamma \alpha \gamma^{-1} = \beta\), so \(\alpha\) and \(\beta\) are conjugate.
\end{proof}
\begin{itemize}
    \item A \textbf{transposition} in \(S_n\) is a 2-cycle
    \item Every element of \(S_n\) can be written as a (not necessarily disjoint) product of transpositions in many ways. However, any two such expressions use the same number, modulo 2, of transpositions
    \item A permutation is \textbf{even} (resp.\ \textbf{odd}) if it is a product of an even (resp.\ odd) number of transpositions; every permutation is exactly one of these. Since an \(r\)-cycle is a product of \(r-1\) transpositions, an \(r\)-cycle is even iff \(r\) is odd
    \item The subset of \(S_n\) consisting of all even permutations is a subgroup of index 2, hence is normal in \(S_n\) by Proposition~\ref{prop:1.8}; this is the \textbf{alternating group} of degree \(n\), denoted \(A_n\)
\end{itemize}
\subsubsection*{Homomorphisms}
\begin{itemize}
    \item Let \(G\) and \(H\) be groups. A \textbf{homomorphism} is a map \(\varphi \colon G \to H\) satisfying \(\varphi(xy) = \varphi(x)\varphi(y)\) for all \(x,y \in G\)
    \item If \(\varphi\) is a homomorphism, then \(\varphi(1) = 1\) and \(\varphi(x^{-1}) = \varphi(x)^{-1}\) for any \(x \in G\)
    \item The \textbf{trivial homomorphism} from \(G\) to \(H\) sends every element of \(G\) to the identity of \(H\)
    \item \(\varphi\) is a \textbf{monomorphism} if injective, an \textbf{epimorphism} if surjective, and an \textbf{isomorphism} if bijective. If \(\varphi\) is an isomorphism, then so is \(\varphi^{-1} : H \to G\)
    \item A homomorphism \(\varphi : G \to G\) is an \textbf{endomorphism} of \(G\); a bijective endomorphism is an \textbf{automorphism}
    \item If there is an isomorphism \(\varphi : G \to H\), we say \(G\) and \(H\) are \textbf{isomorphic} and write \(G \cong H\). Isomorphism is an equivalence relation on groups, so we can speak of the \emph{isomorphism class} of a group.
\end{itemize}
\noindent Standard examples of homomorphisms and isomorphisms:
\begin{itemize}
    \item Let \(G = \langle g \rangle\) and \(H = \langle h \rangle\) be cyclic groups of order \(n\). The map \(\varphi : G \to H\) by \(\varphi(g^a) = h^a\) is an isomorphism. Hence any two finite cyclic groups of the same order are isomorphic. 
    \item Let \(G\) be a group, \(H \leq G\), \(g \in G\). The \textbf{conjugate} of \(H\) by \(g\) is \(gHg^{-1} = \{ghg^{-1} \mid h \in H\}\), and \(gHg^{-1} \leq G\). The map \(\varphi : H \to gHg^{-1}\) by \(\varphi(h) = ghg^{-1}\) is an isomorphism, so conjugate subgroups are isomorphic. (However, isomorphic subgroups need not be conjugate)
    \item Let \(X = \{x_1, \dots, x_n\}\) and let \(S_X\) be the group of permutations of \(X\). The map \(\varphi : S_n \to S_X\) by \(\varphi(\rho)(x_i) = x_{\rho(i)}\) is an isomorphism
    \item Let \(G\) be a group, \(N \trianglelefteq G\). The \textbf{natural map} \(\eta : G \to G/N\) defined by \(\eta(x) = xN\) is an epimorphism
\end{itemize}
\noindent Kernel and image:
\begin{itemize}
    \item If \(\varphi : G \to H\) is a homomorphism, the \textbf{kernel} of \(\varphi\) is \(\ker \varphi = \{g \in G \mid \varphi(g) = 1\}\), and the \textbf{image} is \(\operatorname{im} \varphi = \{\varphi(g) \mid g \in G\}\)
    \item We use \(\varphi(K) = \{\varphi(k) \mid k \in K\}\) for the image of \(K \leq G\)
    \item E.g., if \(N \trianglelefteq G\) and \(\eta : G \to G/N\) is the natural map, then \(\ker \eta = N\) and \(\eta(K) = KN/N\) for any \(K \leq G\). (\(\eta(K) = K/N\) if \(K\) contains \(N\).)
\end{itemize}
\begin{proposition}
\label{prop:1.11}
Let $G$ and $H$ be groups, and let $\varphi : G \to H$ be a homomorphism. Then $\ker \varphi \trianglelefteq G$, and $\varphi(K) \leq H$ for any $K \leq G$.
\end{proposition}
\begin{proof}
    First, \(\varphi(1) = 1\), so \(1 \in \ker \varphi\). If \(x,y \in \ker \varphi\), then \(\varphi(xy^{-1}) = \varphi(x)\varphi(y)^{-1} = 1\), so \(xy^{-1} \in \ker \varphi\). Hence \(\ker \varphi \leq G\). \\
    Let \(x \in \ker \varphi\) and \(g \in G\). Then \(\varphi(gxg^{-1}) = \varphi(g)\varphi(x)\varphi(g^{-1}) = \varphi(g)\varphi(g)^{-1} = 1\). Thus \(gxg^{-1} \in \ker \varphi\), so \(\ker \varphi \trianglelefteq G\). \\
    Now let \(K \leq G\). Since \(1 \in K\), we have \(1 = \varphi(1) \in \varphi(K)\). If \(a,b \in \varphi(K)\), then \(a = \varphi(k_1)\) and \(b = \varphi(k_2)\) for some \(k_1,k_2 \in K\). Hence \(ab^{-1} = \varphi(k_1)\varphi(k_2)^{-1} = \varphi(k_1k_2^{-1}) \in \varphi(K)\). \\
    \(\therefore \varphi(K) \leq H\).
\end{proof}

\begin{namedtheorem}[Fundamental Theorem on Homomorphisms]
\label{thm:fth}
If $G$ and $H$ are groups and $\varphi : G \to H$ is a homomorphism, then there is an isomorphism $\psi : G/K \to \varphi(G)$ such that $\varphi = \psi \circ \eta$, where $K = \ker \varphi$ and $\eta : G \to G/K$ is the natural map; moreover, the map $\psi$ is uniquely determined.
\end{namedtheorem}
\begin{proof}
    Define \(\psi \colon G/K \to \varphi(G)\) by \(\psi(xK) = \varphi(x)\). This is well-defined: if \(xK = yK\), then \(y^{-1}x \in K\), so \(1 = \varphi(y^{-1}x) = \varphi(y)^{-1}\varphi(x)\), and hence \(\varphi(x) = \varphi(y)\). \\
    The map \(\psi\) is a homomorphism, since \(\psi((xK)(yK)) = \psi(xyK) = \varphi(xy) = \varphi(x)\varphi(y) = \psi(xK)\psi(yK)\). \\
    It is surjective by definition of \(\varphi(G)\): if \(h \in \varphi(G)\), then \(h = \varphi(x) = \psi(xK)\) for some \(x \in G\). It is injective since \(\psi(xK) = 1 \Rightarrow \varphi(x) = 1 \Rightarrow x \in K \Rightarrow xK = K\). Hence \(\psi\) is an isomorphism. \\
    For all \(x \in G\), \((\psi \circ \eta)(x) = \psi(xK) = \varphi(x)\), so \(\varphi = \psi \circ \eta\). \\
    Finally, suppose \(\theta \colon G/K \to \varphi(G)\) also satisfies \(\varphi = \theta \circ \eta\). For any \(xK \in G/K\), \(\theta(xK) = \theta(\eta(x)) = \varphi(x) = \psi(xK)\). Thus \(\theta = \psi\), so \(\psi\) is uniquely determined.
\end{proof}

\begin{namedtheorem}[First Isomorphism Theorem]
\label{thm:first-iso}
Let $G$ be a group. If $N \trianglelefteq G$ and $H \leq G$, then $HN/N \cong H/(H \cap N)$.
\end{namedtheorem}
\begin{proof}
    By Proposition~\ref{prop:1.7}, \(HN \leq G\) and \(H \cap N \trianglelefteq H\). Since \(N \trianglelefteq G\) and \(HN \leq G\), we also have \(N \trianglelefteq HN\), so \(HN/N\) is a group. \\
    Let \(\eta \colon G \to G/N\) be the natural map \(\eta(g) = gN\), and let \(\varphi = \eta|_H \colon H \to G/N\). Then \(\varphi\) is a homomorphism, and \(\ker \varphi = \{h \in H \mid hN = N\} = H \cap N\). \\
    For the natural map, \(\eta(K) = KN/N\) for any \(K \leq G\). Taking \(K = H\), we get \(\varphi(H) = \eta(H) = HN/N\), so we may regard \(\varphi\) as a homomorphism \(H \to HN/N\). \\
    By the \nameref{thm:fth}, \(H/\ker \varphi \cong \varphi(H)\). Thus \(H/(H \cap N) \cong HN/N\). \\
    \(\therefore HN/N \cong H/(H \cap N)\).
\end{proof}

\begin{namedtheorem}[Correspondence Theorem]
\label{thm:correspondence}
Let $G$ and $H$ be groups, and let $\varphi : G \to H$ be an epimorphism having kernel $N$. Then there is a bijective correspondence given by $\varphi$ between the set of subgroups of $G$ that contain $N$ and the set of subgroups of $H$. If $K$ is a subgroup of $G$ containing $N$, then this correspondence sends $K$ to $\varphi(K)$; if $L$ is a subgroup of $H$, then the subgroup of $G$ sent to $L$ under this correspondence is $\varphi^{-1}(L) = \{x \in G \mid \varphi(x) \in L\}$. Moreover, if $K_1$ and $K_2$ are subgroups of $G$ containing $N$, then:
\begin{itemize}
    \item $K_2 \leq K_1$ iff $\varphi(K_2) \leq \varphi(K_1)$, and in this case we have $|K_1 : K_2| = |\varphi(K_1) : \varphi(K_2)|$.
    \item $K_2 \trianglelefteq K_1$ iff $\varphi(K_2) \trianglelefteq \varphi(K_1)$, and in this case the map from $K_1/K_2$ to $\varphi(K_1)/\varphi(K_2)$ sending $xK_2$ to $\varphi(x)\varphi(K_2)$ is an isomorphism.
\end{itemize}
\end{namedtheorem}
\begin{proof}
    Let \(\mathcal{S} = \{K \leq G \mid N \leq K\}\) and \(\mathcal{T} = \{L \leq H\}\). For \(K \in \mathcal{S}\), Proposition~\ref{prop:1.11} gives \(\varphi(K) \leq H\). If \(L \in \mathcal{T}\), then \(\varphi^{-1}(L) \leq G\) by the subgroup test, and \(N \leq \varphi^{-1}(L)\), so \(\varphi^{-1}(L) \in \mathcal{S}\). \\
    Define \(\Phi \colon \mathcal{S} \to \mathcal{T}\) by \(\Phi(K) = \varphi(K)\), and define \(\Psi \colon \mathcal{T} \to \mathcal{S}\) by \(\Psi(L) = \varphi^{-1}(L)\). Since \(\varphi\) is surjective, \(\Phi(\Psi(L)) = \varphi(\varphi^{-1}(L)) = L\). If \(K \in \mathcal{S}\), then \(K \subseteq \Psi(\Phi(K))\). Conversely, if \(x \in \Psi(\Phi(K))\), then \(\varphi(x) = \varphi(k)\) for some \(k \in K\), so \(k^{-1}x \in \ker\varphi = N \leq K\), and hence \(x \in K\). Thus \(\Psi(\Phi(K)) = K\). Hence \(\Phi\) and \(\Psi\) are inverse bijections. \\
    Let \(K_1,K_2 \in \mathcal{S}\). If \(K_2 \leq K_1\), then \(\varphi(K_2) \leq \varphi(K_1)\). Conversely, if \(\varphi(K_2) \leq \varphi(K_1)\), then \(K_2 = \Psi(\Phi(K_2)) \leq \Psi(\Phi(K_1)) = K_1\). \\
    Now suppose \(K_2 \leq K_1\). Define \(\alpha \colon K_1/K_2 \to \varphi(K_1)/\varphi(K_2)\) by \(\alpha(xK_2) = \varphi(x)\varphi(K_2)\). If \(xK_2 = yK_2\), then \(y^{-1}x \in K_2\), so \(\varphi(y)^{-1}\varphi(x) \in \varphi(K_2)\), and \(\alpha(xK_2) = \alpha(yK_2)\). Thus \(\alpha\) is well-defined. It is surjective by definition. If \(\alpha(xK_2) = \alpha(yK_2)\), then \(\varphi(y)^{-1}\varphi(x) = \varphi(k)\) for some \(k \in K_2\). Hence \(k^{-1}y^{-1}x \in N \leq K_2\), so \(y^{-1}x \in K_2\), and \(xK_2 = yK_2\). Thus \(\alpha\) is bijective, so \(|K_1:K_2| = |\varphi(K_1):\varphi(K_2)|\). \\
    If \(K_2 \trianglelefteq K_1\), then for \(x \in K_1\), \(k \in K_2\), we have \(\varphi(x)\varphi(k)\varphi(x)^{-1} = \varphi(xkx^{-1}) \in \varphi(K_2)\), so \(\varphi(K_2) \trianglelefteq \varphi(K_1)\). Conversely, suppose \(\varphi(K_2) \trianglelefteq \varphi(K_1)\). If \(x \in K_1\), \(k \in K_2\), then \(\varphi(xkx^{-1}) \in \varphi(K_2)\), so \(xkx^{-1} \in \Psi(\Phi(K_2)) = K_2\). Hence \(K_2 \trianglelefteq K_1\). \\
    In this case, \(\alpha\) is a homomorphism, since \(\alpha((xK_2)(yK_2)) = \alpha(xyK_2) = \varphi(xy)\varphi(K_2) = \varphi(x)\varphi(y)\varphi(K_2) = \alpha(xK_2)\alpha(yK_2)\). Since \(\alpha\) is already bijective, it is an isomorphism. \\
    \(\therefore\) the correspondence has all the stated properties.
\end{proof}
\smallskip
\noindent As a special case of the \nameref{thm:correspondence}, if \(G\) is a group and \(N \trianglelefteq G\), then every subgroup of \(G/N\) has the form \(K/N\) for some subgroup \(K \leq G\) with \(N \leq K\). Here we take \(\varphi\) to be the natural map \(G \to G/N\).

\begin{namedtheorem}[Second Isomorphism Theorem]
\label{thm:second-iso}
Let $H$ and $K$ be normal subgroups of a group $G$. If $H$ contains $K$, then $G/H \cong (G/K)/(H/K)$.
\end{namedtheorem}
\begin{proof}
    Let \(\eta \colon G \to G/K\) be the natural map. Then \(\ker \eta = K\). Since \(K \leq H \leq G\), the \nameref{thm:correspondence} applies to \(H \leq G\). \\
    Since \(H \trianglelefteq G\), applying the normality and quotient-isomorphism parts with \(K_1 = G\) and \(K_2 = H\) gives \(\eta(H) \trianglelefteq \eta(G)\) and \(G/H \cong \eta(G)/\eta(H)\). \\
    Now \(\eta(G) = G/K\) and \(\eta(H) = H/K\). Hence \(G/H \cong (G/K)/(H/K)\). \\
    \(\therefore G/H \cong (G/K)/(H/K)\).
\end{proof}

\subsection*{Exercises}

\begin{exercise}
\label{ex:1.1}
Prove, or complete the sketched proof of, each result in this section.
\end{exercise}
\begin{solution}
The sketched proofs are completed above in:
\begin{itemize}
    \item Proposition~\ref{prop:1.1}, \nameref{thm:lagrange}, Proposition~\ref{prop:1.2}, and Proposition~\ref{prop:1.3}
    \item Theorem~\ref{thm:1.4}, Theorem~\ref{thm:1.5}, and Theorem~\ref{thm:1.6}
    \item Proposition~\ref{prop:1.7}, Proposition~\ref{prop:1.8}, Theorem~\ref{thm:1.9}, Proposition~\ref{prop:1.10}, and Proposition~\ref{prop:1.11}
    \item \nameref{thm:fth}, \nameref{thm:first-iso}, \nameref{thm:correspondence}, and \nameref{thm:second-iso}
\end{itemize}
\end{solution}

\begin{exercise}
\label{ex:1.2}
We say that a group $G$ has \emph{exponent} $e$ if $e$ is the smallest positive integer such that $x^e = 1$ for every $x \in G$. Show that if $G$ has exponent 2, then $G$ is abelian. For what integers $e$ is a group having exponent $e$ necessarily abelian?
\end{exercise}
\begin{solution}
First suppose that \(G\) has exponent \(2\). Let \(x,y \in G\). Then \(x^2 = y^2 = (xy)^2 = 1\), so \(xyxy = 1\). Multiplying on the left by \(x\) and on the right by \(y\), and using \(x^2 = y^2 = 1\), gives \(yx = xy\). Since \(x,y\) were arbitrary, \(G\) is abelian. \\
If \(e = 1\), then \(x = 1\) for all \(x \in G\), so \(G = \{1\}\), which is abelian. We claim that these are the only possibilities. \\
Let \(e \geq 3\). If \(e\) is even, take the dihedral group
\[
    D_e = \langle r,s \mid r^e = s^2 = 1,\ srs^{-1} = r^{-1}\rangle.
\]
This group is non-abelian, since \(srs^{-1} = r^{-1} \ne r\). Its rotations have order dividing \(e\), and its reflections have order \(2\). Since \(e\) is even, the exponent is \(\operatorname{lcm}(e,2) = e\). \\
Now suppose that \(e \geq 3\) is odd. Let
\[
    U_e =
    \left\{
    M(a,b,c) =
    \begin{pmatrix}
        1 & a & c \\
        0 & 1 & b \\
        0 & 0 & 1
    \end{pmatrix}
    \ \middle|\ a,b,c \in \Z/e\Z
    \right\},
\]
under matrix multiplication, with entries computed in \(\Z/e\Z\). Direct multiplication gives
\[
    M(a,b,c)M(a',b',c') = M(a+a', b+b', c+c'+ab').
\]
Hence \(U_e\) is closed under multiplication. The identity is \(M(0,0,0)\), and
\[
    M(a,b,c)^{-1} = M(-a,-b,ab-c).
\]
Associativity follows from associativity of matrix multiplication, so \(U_e\) is a group. It is non-abelian, since
\[
    M(1,0,0)M(0,1,0) = M(1,1,1),
    \qquad
    M(0,1,0)M(1,0,0) = M(1,1,0).
\]
For \(n \in \N\), induction using the multiplication formula gives
\[
    M(a,b,c)^n = M\left(na, nb, nc+\binom{n}{2}ab\right).
\]
Taking \(n=e\), we get
\[
    M(a,b,c)^e = M\left(0,0,\binom{e}{2}ab\right).
\]
Since \(e\) is odd, \(\binom{e}{2} = e(e-1)/2\) is divisible by \(e\), so \(M(a,b,c)^e = M(0,0,0)\) for every \(M(a,b,c) \in U_e\). Thus every element has order dividing \(e\). However \(M(1,0,0)^n = M(n,0,0)\), so \(M(1,0,0)\) has order \(e\). Hence \(U_e\) is a non-abelian group of exponent \(e\). \\
\(\therefore\) a group of exponent \(e\) is necessarily abelian iff \(e = 1\) or \(e = 2\).
\end{solution}

\begin{exercise}
\label{ex:1.3}
Let $G$ be a finite group, and suppose that the map $\varphi : G \to G$ defined by $\varphi(x) = x^3$ for $x \in G$ is a homomorphism. Show that if 3 does not divide $|G|$, then $G$ must be abelian. (See [2] for a generalization.)
\end{exercise}
\begin{solution}
Let \(x,y \in G\). Since \(\varphi\) is a homomorphism, we have
\[
    x^3y^3 = \varphi(x)\varphi(y) = \varphi(xy) = (xy)^3 = xyxyxy.
\]
Cancelling \(x\) on the left and \(y\) on the right gives
\[
    (yx)^2 = x^2y^2.
\]
Swapping \(x\) and \(y\), we also have
\[
    (xy)^2 = y^2x^2.
\]
Using this in the original equation,
\[
    x^3y^3 = (xy)^3 = (xy)(xy)^2 = xy(y^2x^2) = xy^3x^2.
\]
Cancelling \(x\) on the left gives
\[
    y^3x^2 = x^2y^3.
\]
Hence every cube commutes with every square. \\
Now let \(z \in \ker \varphi\). Then \(z^3 = 1\), so \(o(z) \mid 3\). Then we also have \(o(z) \mid |G|\). Since \(3 \nmid |G|\), we must have \(o(z) = 1\), and hence \(z = 1\). Thus \(\ker \varphi = \{1\}\), so \(\varphi\) is injective. Since \(G\) is finite, \(\varphi\) is also surjective. Therefore every element of \(G\) is a cube. \\
Let \(g \in G\). Since every element of \(G\) is a cube, write \(g = y^3\) for some \(y \in G\). Then for every \(x \in G\), the relation above gives
\[
    gx^2 = y^3x^2 = x^2y^3 = x^2g.
\]
Thus every square is central. In particular, \(x^2\) and \(y^2\) commute for all \(x,y \in G\). Since \((xy)^2 = y^2x^2\), it follows that
\[
    (xy)^2 = x^2y^2.
\]
Hence \(xyxy = xxyy\). Cancelling \(x\) on the left and \(y\) on the right gives \(yx = xy\). Since \(x,y\) were arbitrary, \(G\) is abelian.
\end{solution}

\begin{exercise}
\label{ex:1.4}
Let $g$ be an element of a group $G$, and suppose that $|G| = mn$ where $m$ and $n$ are coprime. Show that there are unique elements $x$ and $y$ of $G$ such that $xy = g = yx$ and $x^m = 1 = y^n$. (In the case where $m$ is a power of some prime $p$, we call $x$ the $p$-part of $g$ and $y$ the $p'$-part of $g$; more generally, if $\pi$ is a set of primes which includes all prime divisors of $m$ and no prime divisors of $n$, then $x$ and $y$ are called the $\pi$-part and $\pi'$-part, respectively, of $g$.)
\end{exercise}
\begin{solution}
Since \(\gcd(m,n) = 1\), choose \(a,b \in \Z\) such that \(am+bn = 1\). Since \(\langle g \rangle \leq G\), \nameref{thm:lagrange} gives \(o(g) = |\langle g \rangle| \mid |G| = mn\), so \(g^{mn} = 1\). Define
\[
    x = g^{bn}, \qquad y = g^{am}.
\]
Then \(x\) and \(y\) commute, since both are powers of \(g\). Moreover,
\[
    xy = g^{bn}g^{am} = g^{am+bn} = g,
\]
and similarly \(yx = g\). Also,
\[
    x^m = g^{bmn} = (g^{mn})^b = 1,
    \qquad
    y^n = g^{amn} = (g^{mn})^a = 1.
\]
Thus such elements \(x,y \in G\) exist. \\
We now show that they are unique. Suppose that \(x',y' \in G\) also satisfy
\[
    x'y' = g = y'x',
    \qquad
    (x')^m = 1,
    \qquad
    (y')^n = 1.
\]
Since \(x'\) and \(y'\) commute, we may take powers of \(x'y'\) termwise. Thus
\[
    g^{bn} = (x'y')^{bn} = (x')^{bn}(y')^{bn}.
\]
But \((y')^{bn} = ((y')^n)^b = 1\), so
\[
    g^{bn} = (x')^{bn}.
\]
Since \(bn = 1-am\), we get
\[
    (x')^{bn} = (x')^{1-am} = x'((x')^m)^{-a} = x'.
\]
Therefore \(x' = g^{bn} = x\). Similarly,
\[
    g^{am} = (x'y')^{am} = (x')^{am}(y')^{am} = (y')^{am},
\]
and since \(am = 1-bn\),
\[
    (y')^{am} = (y')^{1-bn} = y'((y')^n)^{-b} = y'.
\]
Therefore \(y' = g^{am} = y\). Hence the required elements \(x\) and \(y\) exist and are unique. \\
\end{solution}

\begin{exercise}
\label{ex:1.5}
Let $r, s$, and $t$ be positive integers greater than 1. Show that there is a finite group $G$ having elements $x$ and $y$ such that $x$ has order $r$, $y$ has order $s$, and $xy$ has order $t$.
\end{exercise}
\begin{solution}
We use one standard fact about a particular group given by generators and relations. For \(a,b,c > 1\), the triangle group \(\Delta(a,b,c)\) is the abstract group
\[
    \Delta(a,b,c)
    =
    \langle U,V \mid U^a = 1,\ V^b = 1,\ (UV)^c = 1\rangle
\]
generated by two formal elements \(U\) and \(V\), subject to the relations \(U^a=1\), \(V^b=1\), and \((UV)^c=1\). The word ``triangle'' is only the standard name for this presentation; no geometry is needed here. Equivalently, \(\Delta(a,b,c)\) is the universal group with these three relations: whenever a group \(H\) has elements \(\alpha,\beta \in H\) satisfying
\[
    \alpha^a = 1,
    \qquad
    \beta^b = 1,
    \qquad
    (\alpha\beta)^c = 1,
\]
there is a homomorphism \(\Delta(a,b,c) \to H\) sending \(U\) to \(\alpha\) and \(V\) to \(\beta\). \\
The standard fact we use is that \(U,V,UV\) have exact orders \(a,b,c\), respectively, in \(\Delta(a,b,c)\), and that \(\Delta(a,b,c)\) is residually finite. Here residually finite means that if \(g \ne 1\), then there is a homomorphism from \(\Delta(a,b,c)\) to some finite group in which the image of \(g\) is not \(1\). \\
Now apply this fact with \(a=r\), \(b=s\), and \(c=t\). Put
\[
    \Gamma =
    \langle X,Y \mid X^r = 1,\ Y^s = 1,\ (XY)^t = 1\rangle .
\]
In \(\Gamma\), the elements \(X,Y,XY\) have exact orders \(r,s,t\). Hence every element of the finite set
\[
    \mathcal S
    =
    \{X^i \mid 1 \leq i < r\}
    \cup
    \{Y^j \mid 1 \leq j < s\}
    \cup
    \{(XY)^k \mid 1 \leq k < t\}
\]
is non-identity. By residual finiteness, for each \(g \in \mathcal S\), there is a finite group \(Q_g\) and a homomorphism
\[
    \varphi_g \colon \Gamma \to Q_g
\]
such that \(\varphi_g(g) \ne 1\). Now define
\[
    Q = \prod_{g \in \mathcal S} Q_g,
    \qquad
    \Phi \colon \Gamma \to Q,
    \qquad
    \Phi(h) = \bigl(\varphi_g(h)\bigr)_{g \in \mathcal S}.
\]
Let \(G = \Phi(\Gamma) \leq Q\), and set
\[
    x = \Phi(X),
    \qquad
    y = \Phi(Y).
\]
Since \(Q\) is finite, \(G\) is finite. Also, since \(X^r = 1\), \(Y^s = 1\), and \((XY)^t = 1\) in \(\Gamma\), we have
\[
    x^r = 1,
    \qquad
    y^s = 1,
    \qquad
    (xy)^t = 1.
\]
Thus \(o(x) \mid r\), \(o(y) \mid s\), and \(o(xy) \mid t\). \\
It remains to show that no smaller positive powers become trivial in this finite quotient. Suppose \(1 \leq i < r\). Since \(X^i \in \mathcal S\), the \(X^i\)-coordinate of \(\Phi(X^i)\) is \(\varphi_{X^i}(X^i)\), which is not \(1\). Hence \(\Phi(X^i) \ne 1\), so \(x^i \ne 1\). Therefore \(o(x) = r\). \\
Similarly, if \(1 \leq j < s\), then the \(Y^j\)-coordinate of \(\Phi(Y^j)\) is \(\varphi_{Y^j}(Y^j) \ne 1\), so \(y^j \ne 1\). Hence \(o(y) = s\). Finally, if \(1 \leq k < t\), then the \((XY)^k\)-coordinate of \(\Phi((XY)^k)\) is \(\varphi_{(XY)^k}((XY)^k) \ne 1\), so \((xy)^k \ne 1\). Hence \(o(xy) = t\). \\
\(\therefore\) the finite group \(G\) contains elements \(x,y\) such that \(o(x)=r\), \(o(y)=s\), and \(o(xy)=t\).
\end{solution}

\begin{exercise}
\label{ex:1.6}
Let $X$ and $Y$ be subsets of a group $G$. Are $\langle X \rangle \cap \langle Y \rangle$ and $\langle X \cap Y \rangle$ necessarily equal? Are $\langle \langle X \rangle \cup \langle Y \rangle \rangle$ and $\langle X \cup Y \rangle$ necessarily equal?
\end{exercise}
\begin{solution}
The first pair need not be equal. Let \(G = C_3 = \langle g \rangle\), \(X = \{1,g\}\), and \(Y = \{1,g^2\}\). Then \(\langle X\rangle = G\), since \(g \in X\), and \(\langle Y\rangle = G\), since \(g^2\) also generates \(C_3\). Hence
\[
    \langle X\rangle \cap \langle Y\rangle = G.
\]
However, \(X \cap Y = \{1\}\), so
\[
    \langle X \cap Y\rangle = \{1\}.
\]
Thus \(\langle X\rangle \cap \langle Y\rangle\) and \(\langle X \cap Y\rangle\) are not necessarily equal. \\
The second pair is always equal. Since \(X \subseteq \langle X\rangle\) and \(Y \subseteq \langle Y\rangle\), we have
\[
    X \cup Y \subseteq \langle X\rangle \cup \langle Y\rangle.
\]
Therefore
\[
    \langle X \cup Y\rangle \leq \langle \langle X\rangle \cup \langle Y\rangle\rangle.
\]
Conversely, since \(X \subseteq X \cup Y\) and \(Y \subseteq X \cup Y\), we have
\[
    \langle X\rangle \leq \langle X \cup Y\rangle,
    \qquad
    \langle Y\rangle \leq \langle X \cup Y\rangle.
\]
Hence
\[
    \langle X\rangle \cup \langle Y\rangle \subseteq \langle X \cup Y\rangle.
\]
Since \(\langle X \cup Y\rangle\) is a subgroup containing \(\langle X\rangle \cup \langle Y\rangle\), it contains the subgroup generated by that union, so
\[
    \langle \langle X\rangle \cup \langle Y\rangle\rangle \leq \langle X \cup Y\rangle.
\]
\(\therefore \langle \langle X\rangle \cup \langle Y\rangle\rangle = \langle X \cup Y\rangle\).
\end{solution}

\begin{exercise}
\label{ex:1.7}
Let $G$ be a finite group and let $H \leq G$. Show that there is a subset $T$ of $G$ which is simultaneously a left transversal for $H$ and a right transversal for $H$.
\end{exercise}
\begin{solution}
We use the following elementary matching fact. Suppose \(A_1,\dots,A_n\) are finite sets of possible choices. If, for every \(I \subseteq \{1,\dots,n\}\),
\[
    \left|\bigcup_{i \in I} A_i\right| \geq |I|,
\]
then we can choose distinct representatives \(a_i \in A_i\). \\
Let
\[
    L_1,\dots,L_n
\]
be the distinct left cosets of \(H\), and let
\[
    R_1,\dots,R_n
\]
be the distinct right cosets of \(H\). There are the same number \(n\) of each, since both left and right cosets partition \(G\) into subsets of size \(|H|\). For each \(i\), define
\[
    A_i = \{j \in \{1,\dots,n\} \mid L_i \cap R_j \ne \emptyset\}.
\]
Thus \(A_i\) is the set of right cosets which the left coset \(L_i\) meets. We check the matching condition. Choose any \(m\) left cosets, say those indexed by \(I\), where \(|I|=m\). Suppose these \(m\) left cosets meet \(q\) right cosets. Since distinct left cosets are disjoint,
\[
    \left|\bigcup_{i \in I} L_i\right| = m|H|.
\]
This union is contained in the union of the \(q\) right cosets it meets, whose total size is \(q|H|\). Hence
\[
    m|H| \leq q|H|,
\]
so \(m \leq q\). Therefore any \(m\) left cosets meet at least \(m\) right cosets. In terms of the sets \(A_i\), this says exactly that
\[
    \left|\bigcup_{i \in I} A_i\right| \geq |I|.
\]
By the matching fact, choose distinct \(j_1,\dots,j_n\) such that \(j_i \in A_i\) for every \(i\). Since \(j_i \in A_i\),
\[
    L_i \cap R_{j_i} \ne \emptyset,
\]
so choose
\[
    t_i \in L_i \cap R_{j_i}.
\]
Set
\[
    T = \{t_1,\dots,t_n\}.
\]
Because \(t_i \in L_i\), the set \(T\) contains one element from each left coset. Since the \(L_i\) are disjoint, it contains exactly one element from each left coset, so \(T\) is a left transversal for \(H\). \\
Because the indices \(j_1,\dots,j_n\) are distinct, the right cosets \(R_{j_1},\dots,R_{j_n}\) are all distinct. There are \(n\) of them, and there are only \(n\) right cosets in total, so they are all the right cosets. Since \(t_i \in R_{j_i}\), the set \(T\) also contains exactly one element from each right coset. Hence \(T\) is a right transversal for \(H\). \\
\(\therefore\) there exists a subset \(T \subseteq G\) which is simultaneously a left transversal and a right transversal for \(H\).
\end{solution}

\begin{exercise}
\label{ex:1.8}
Suppose that $\mathcal{C}$ is a family of subsets of a group $G$ which forms a partition of $G$, and suppose further that $gC \in \mathcal{C}$ for any $g \in G$ and $C \in \mathcal{C}$. (Recall that a \emph{partition} of a set $S$ is a collection $\mathcal{S}$ of subsets of $S$ with the property that every element of $S$ lies in exactly one member of $\mathcal{S}$.) Show that $\mathcal{C}$ is the set of cosets of some subgroup of $G$.
\end{exercise}
\begin{solution}
Since \(\mathcal C\) is a partition of \(G\), there is a unique member \(H \in \mathcal C\) such that \(1 \in H\). We first show that \(H \leq G\). \\
Let \(h \in H\). By the hypothesis, \(hH \in \mathcal C\). Also, since \(1 \in H\), we have \(h = h1 \in hH\). Thus both \(H\) and \(hH\) are members of \(\mathcal C\) containing \(h\). Since \(\mathcal C\) is a partition, we must have
\[
    hH = H.
\]
Hence for any \(h,k \in H\), we have \(hk \in hH = H\), so \(H\) is closed under multiplication. Also, since \(1 \in H = hH\), there exists \(u \in H\) such that \(hu = 1\). Thus \(u = h^{-1}\), so \(h^{-1} \in H\). Therefore \(H \leq G\). \\
Now let \(C \in \mathcal C\). Choose \(g \in C\). By the hypothesis, \(gH \in \mathcal C\), and since \(1 \in H\), we have \(g \in gH\). Thus both \(C\) and \(gH\) are members of \(\mathcal C\) containing \(g\). Since \(\mathcal C\) is a partition, \(C = gH\). Therefore every member of \(\mathcal C\) is a left coset of \(H\). Conversely, for every \(g \in G\), the set \(gH\) belongs to \(\mathcal C\) by the hypothesis. \\
\(\therefore \mathcal C = \{gH \mid g \in G\}\), the set of left cosets of the subgroup \(H\).
\end{solution}

\begin{exercise}
\label{ex:1.9}
Suppose that $\mathcal{C}$ is a family of subsets of a group $G$ which forms a partition of $G$, and suppose further that $XY \in \mathcal{C}$ for any $X, Y \in \mathcal{C}$. Show that exactly one of the sets belonging to $\mathcal{C}$ is a subgroup of $G$, that this subgroup is normal in $G$, and that $\mathcal{C}$ consists of its cosets.
\end{exercise}
\begin{solution}
Since \(\mathcal C\) is a partition of \(G\), there is a unique member \(H \in \mathcal C\) such that \(1 \in H\). We first show that \(H \leq G\). \\
By the hypothesis, \(HH \in \mathcal C\). Since \(1 \in H\), we have \(1 = 1 \cdot 1 \in HH\). Thus \(H\) and \(HH\) are both members of \(\mathcal C\) containing \(1\), so \(HH = H\). Hence \(H\) is closed under multiplication. \\
Now let \(h \in H\). Let \(K \in \mathcal C\) be the unique member containing \(h^{-1}\). By the hypothesis, \(HK \in \mathcal C\). Since \(h \in H\) and \(h^{-1} \in K\), we have \(1 = hh^{-1} \in HK\). Thus \(HK\) and \(H\) are both members of \(\mathcal C\) containing \(1\), so \(HK = H\). Since \(1 \in H\) and \(h^{-1} \in K\), it follows that \(h^{-1} = 1h^{-1} \in HK = H\). Therefore \(H \leq G\). \\
We now show that every member of \(\mathcal C\) is a left coset of \(H\). Let \(X \in \mathcal C\), and choose \(x \in X\). By the hypothesis, \(XH \in \mathcal C\). Since \(1 \in H\), we have \(x = x1 \in XH\). Thus \(XH\) and \(X\) are both members of \(\mathcal C\) containing \(x\), so \(XH = X\). Hence \(xH \subseteq XH = X\). \\
For the reverse inclusion, let \(X' \in \mathcal C\) be the unique member containing \(x^{-1}\). By the hypothesis, \(X'X \in \mathcal C\). Since \(x^{-1}x = 1 \in X'X\), we have \(X'X = H\). If \(y \in X\), then \(x^{-1}y \in X'X = H\), and hence
\[
    y = x(x^{-1}y) \in xH.
\]
Thus \(X \subseteq xH\). Therefore \(X = xH\), and so every member of \(\mathcal C\) is a left coset of \(H\). Conversely, if \(g \in G\), and \(X \in \mathcal C\) contains \(g\), then \(X = gH\). Hence
\[
    \mathcal C = \{gH \mid g \in G\}.
\]
The subgroup \(H\) is the only member of \(\mathcal C\) which is a subgroup, since it is the only member containing \(1\). \\
It remains to show that \(H\) is normal. Let \(g \in G\). Since \(gH,g^{-1}H \in \mathcal C\), the product \((gH)(g^{-1}H)\) is a member of \(\mathcal C\). It contains \(1\), since \(g \in gH\), \(g^{-1} \in g^{-1}H\), and \(gg^{-1}=1\). Thus
\[
    (gH)(g^{-1}H)=H.
\]
Now let \(h \in H\). Since \(gh \in gH\) and \(g^{-1} \in g^{-1}H\), we have
\[
    ghg^{-1} \in (gH)(g^{-1}H)=H.
\]
Thus \(gHg^{-1} \subseteq H\). Replacing \(g\) by \(g^{-1}\), we also have \(g^{-1}Hg \subseteq H\), and conjugating this inclusion by \(g\) gives \(H \subseteq gHg^{-1}\). Hence \(gHg^{-1}=H\) for all \(g \in G\). Therefore \(H \trianglelefteq G\). \\
\(\therefore\) exactly one member of \(\mathcal C\) is a subgroup, this subgroup is normal in \(G\), and \(\mathcal C\) consists of its cosets.
\end{solution}

\begin{exercise}
\label{ex:1.10}
Prove the following generalization of Proposition~\ref{prop:1.8}: If $G$ is a finite group and $H \leq G$ is such that $|G : H|$ is equal to the smallest prime divisor of $|G|$, then $H \trianglelefteq G$.
\end{exercise}
\begin{solution}
Let \(p = |G:H|\). By hypothesis, \(p\) is the smallest prime divisor of \(|G|\). Let
\[
    \Omega = \{xH \mid x \in G\}
\]
be the set of left cosets of \(H\). Then \(|\Omega|=p\). Define
\[
    \varphi \colon G \to S_\Omega
\]
by
\[
    \varphi(g)(xH) = gxH
\]
for all \(g,x \in G\). This is well-defined: if \(xH = yH\), then \(gxH = gyH\). For each \(g \in G\), the map \(\varphi(g)\) is a permutation of \(\Omega\), with inverse \(\varphi(g^{-1})\). Also,
\[
    \varphi(g_1g_2)(xH) = g_1g_2xH
    =
    \varphi(g_1)(g_2xH)
    =
    \bigl(\varphi(g_1)\varphi(g_2)\bigr)(xH),
\]
so \(\varphi\) is a homomorphism. \\
Let \(K = \ker \varphi\). Then \(K \trianglelefteq G\). We claim that \(K \leq H\). If \(k \in K\), then \(\varphi(k)\) fixes every left coset, so in particular
\[
    H = \varphi(k)(H) = kH.
\]
Thus \(k \in H\), and hence \(K \leq H\). \\
By the \nameref{thm:fth}, \(G/K \cong \operatorname{im}\varphi \leq S_\Omega\). Hence
\[
    |G:K| = |\operatorname{im}\varphi| \mid |S_\Omega| = p!.
\]
Also, by \nameref{thm:lagrange}, \(|G:K| \mid |G|\). Therefore every prime divisor of \(|G:K|\) is a prime divisor of \(|G|\), and hence is at least \(p\). On the other hand, since \(|G:K| \mid p!\), every prime divisor of \(|G:K|\) is at most \(p\). Thus the only possible prime divisor of \(|G:K|\) is \(p\). \\
Since \(K \leq H \leq G\), we have
\[
    |G:K| = |G:H|\,|H:K| = p|H:K|.
\]
Thus \(p \mid |G:K|\). Since the only possible prime divisor of \(|G:K|\) is \(p\), we may write \(|G:K|=p^a\) for some \(a \geq 1\). But \(|G:K| \mid p!\), and \(p^2 \nmid p!\), so \(a=1\). Hence
\[
    |G:K|=p.
\]
It follows that \(p|H:K|=p\), so \(|H:K|=1\). Therefore \(H=K\). Since \(K \trianglelefteq G\), we conclude that \(H \trianglelefteq G\).
\end{solution}

\subsection*{Further Exercises}

\noindent If $K \trianglelefteq H \leq G$, then $H/K$ is called a \emph{section} of $G$. We say that two sections $H_1/K_1$ and $H_2/K_2$ of $G$ are \emph{incident} if every coset of $K_1$ in $H_1$ has non-empty intersection with exactly one coset of $K_2$ in $H_2$, and vice versa. (In other words, two sections are incident if the relation of non-empty intersection gives a bijective correspondence between their elements.)

\begin{exercise}
\label{ex:1.11}
Show that incident sections are isomorphic.
\end{exercise}
\begin{solution}
\end{solution}

\begin{exercise}
\label{ex:1.12}
(cont.) Suppose that $N \trianglelefteq G$ and $H \leq G$. Show that $HN/N$ and $H/H \cap N$ are incident. (Exercises~\ref{ex:1.11} and~\ref{ex:1.12} provide an alternate proof of the \nameref{thm:first-iso}.)
\end{exercise}
\begin{solution}
\end{solution}

\noindent If $L/M$ is a section of $G$ and $H \leq G$, then the \emph{projection} of $H$ on $L/M$ is the subset of $L/M$ consisting of those cosets of $M$ in $L$ which contain elements of $H$.

\begin{exercise}
\label{ex:1.13}
(cont.) Show that the projection of $H$ on $L/M$ is the subgroup $(L \cap H)M/M$ of $L/M$.
\end{exercise}
\begin{solution}
\end{solution}

\noindent Let $H_1/K_1$ and $H_2/K_2$ be sections of a group $G$.

\begin{exercise}
\label{ex:1.14}
(cont.) Show that the projection of $K_2$ on $H_1/K_1$ is a normal subgroup of the projection of $H_2$ on $H_1/K_1$. The quotient group obtained thereby is called the projection of $H_2/K_2$ on $H_1/K_1$.
\end{exercise}
\begin{solution}
\end{solution}

\begin{namedtheorem}[Third Isomorphism Theorem]
\label{thm:third-iso}
Let $H_1, H_2 \leq G$, let $K_1 \trianglelefteq H_1$, and let $K_2 \trianglelefteq H_2$. Then
\[(H_1 \cap H_2)K_1/(H_1 \cap K_2)K_1 \cong (H_1 \cap H_2)K_2/(K_1 \cap H_2)K_2.\]
\end{namedtheorem}
\begin{proof}
\end{proof}
\smallskip
\noindent (This result is also called the fourth isomorphism theorem, or Zassenhaus' lemma --- after its discoverer, who proved it as a student at the age of 21 --- or even the butterfly lemma. The last name refers to the shape of the diagram showing the inclusion relations between the many subgroups involved.)

\begin{exercise}
\label{ex:1.15}
(cont.) Show that the projection of $H_1/K_1$ on $H_2/K_2$ and the projection of $H_2/K_2$ on $H_1/K_1$ are incident. Deduce the \nameref{thm:third-iso}.
\end{exercise}
\begin{solution}
\end{solution}

\end{document}
